% ======================================================================
% SECTION 1: INTRODUCTION
% Five thematic subsections: FDA RTCT, GBM clinical context, current
% robotic neurosurgery baseline, on-premises LLM thesis, and the
% transition to the 60-second resection.
% Location: 2030-gbm-1min/paper/full-paper/sections/introduction.tex
% ======================================================================

\subsection{FDA Real-Time Clinical Trials Announcement}
\label{sec:intro-fda}

On 28 April 2026 the United States Food and Drug Administration (FDA)
announced that real-time clinical trials (RTCT) are entering the
proof-of-concept phase \cite{fda2026realtime}. The announcement names
two oncology pharmacology pilots: AstraZeneca's TRAVERSE in mantle cell
lymphoma and Amgen's STREAM-SCLC in small cell lung carcinoma. FDA
Commissioner Marty Makary framed the RTCT program as a structural
upgrade to the clinical-trial data pipeline; FDA Chief AI Officer
Jeremy Walsh tied the program to the agency's broader software and AI
modernization roadmap.

The 60-second glioblastoma robotic surgery target studied in this
paper is positioned as a natural extension of the RTCT vision into
the surgical theater, not a critique of the RTCT program. The two
proof-of-concept trials cover pharmacology only; the 60-second
glioblastoma resection extends real-time signal sharing into the
surgical theater through the on-premises repository based LLM thesis
described in \S\ref{sec:intro-thesis}. The United States needs to
remain Number 1 in the world regarding patient safety, efficacy, and
speed benefits in oncological robotic surgeries in clinical trials;
the on-premises LLM thesis is one route to that goal, and the
present work provides a reproducible computational checking step that
a downstream NeuroSpeed-class hardware vendor or trial sponsor can
adopt without having to re-author the simulation pipeline. Prior
author work positioned a parallel 1-hour single-arm pharmacology-plus-
robotics framing under the same FDA roadmap
\cite{kawchak_2026_19994945}; the present work compresses the
procedure dimension by 60x to surface 4-arm signal-sharing dynamics
that the 1-hour scope did not require.

\subsection{Glioblastoma Clinical Context}
\label{sec:intro-gbm}

IDH-wildtype glioblastoma multiforme (WHO grade 4) remains a uniformly
lethal primary brain malignancy. Median survival under standard-of-care
radiotherapy plus concomitant and adjuvant temozolomide is 12 to 15
months \cite{Stupp2005Glioblastoma}. The extent-of-resection threshold
for newly diagnosed glioblastoma is 78 percent of contrast-enhancing
volume on T1 magnetic resonance imaging \cite{Sanai2011GBMResection};
5-aminolevulinic acid fluorescence-guided resection (5-ALA) increases
the rate of complete resection of contrast-enhancing tumor compared
to white-light microsurgery \cite{Stummer20055ALA}. The synthetic
patient PAT-GBM-0001 used throughout this paper is a 62-year-old
female with a 4.2 cm right frontal IDH-wildtype glioblastoma, MGMT
methylated, 1.8 cm to the motor strip on the non-dominant side; the
maximal-safe gross-total resection target follows the
Stupp-Sanai-Stummer framework. The patient identity, disease
parameters, and pre-operative imaging are inherited verbatim from the
upstream physical-ai-oncology-trials instruction directory
\cite{one-minute-variant-instructions}.

Current operative duration for a single-arm stereotactic-guided right
frontal craniotomy with maximal safe resection ranges between 4 and
8 hours of in-room time. A 60-second target is a 240x to 480x
compression that is not feasible with any current commercial robot.
This contrast frames the rest of the paper as a checking step the
industry has not yet taken. The author's prior glioblastoma work
spans an autonomous 10-year meta-analysis effort
\cite{kawchak_2025_15549831}, drug-synergy machine learning
\cite{kawchak_2025_17614396}, peer-review acceleration of patient
matching \cite{kawchak_2025_17774560}, and the parallel 1-hour
robotic framing \cite{kawchak_2026_19994945}; the present 60-second
target is the next computational checkpoint along that progression.

\subsection{Current Robotic Neurosurgery State of the Art}
\label{sec:intro-baseline}

The state of the art for stereotactic neurosurgical robotics in 2026
is the Medtronic ROSA ONE Brain v3.0 \cite{rosa-one-brain}, used
clinically for stereotactic biopsy, deep-brain-stimulation lead
placement, and stereoelectroencephalography
\cite{Lefranc2014ROSAStereotaxy, Kalakoti2019RoboticsReview}. ROSA is
single arm and is not designed for high-throughput soft-tissue
resection. Table \ref{tab:robot-comparison} summarizes the side-by-
side specification against the hypothetical 2030 Medtronic NeuroSpeed
1.0 across the parameters that gate a 60-second resection. NeuroSpeed
1.0 is paper-only at the time of writing; the hardware does not exist.

\begin{table}[H]
\centering
\caption{Specification comparison: Medtronic ROSA ONE Brain v3.0
(current state of the art, single arm) vs hypothetical 2030 Medtronic
NeuroSpeed 1.0 (4 cooperating arms). All gaps are 5x to 200x.}
\label{tab:robot-comparison}
\begin{tabular}{>{\raggedright\arraybackslash}p{4.5cm}
                >{\raggedright\arraybackslash}p{3.3cm}
                >{\raggedright\arraybackslash}p{3.3cm}
                >{\raggedright\arraybackslash}p{2.0cm}}
\toprule
\textbf{Parameter} & \textbf{ROSA ONE Brain v3.0} & \textbf{NeuroSpeed 1.0 (2030)} & \textbf{Gap} \\
\midrule
Tip linear velocity & 50 mm/s & 1,000 mm/s & 20x \\
Tip linear acceleration & 200 mm/s squared & 10,000 mm/s squared & 50x \\
Joint angular velocity & 180 deg/s & 360 deg/s & 2x \\
Positioning accuracy at speed & 0.5 mm RMS & 0.1 mm RMS & 5x \\
Force resolution & 0.01 N & 0.001 N & 10x \\
E-stop latency & 50 ms & 5 ms & 10x \\
Tissue removal peak rate & 4 mm cubed/s per arm & 800 mm cubed/s per arm & 200x \\
Per-arm tip force limit & 15.0 N & 5.0 N & tighter \\
Cumulative 4-arm tip force & not specified & 12.0 N envelope & new \\
Cooperating arms & 1 & 4 (7 DOF each) & 4x \\
Peak duty cycle & 4 to 8 hours & 5 minute peak (LN2 cooled) & by design \\
\bottomrule
\end{tabular}
\end{table}

NeuroSpeed 1.0 closes the gaps through liquid-nitrogen cooling, a
hybrid ultrasonic plus waterjet plus pulsed plasma cutting head on
arm 1, and a deterministic 1 kHz heartbeat broadcast bus that
coordinates the 4 arms. The regulatory floor is set by IEC 80601-2-77
(per-arm 5.0 N tip and 1.0 N lateral force limits, 5 ms E-stop
budget) \cite{iec-80601-2-77} and by IEC 62304 (safety-critical
software lifecycle for the 1 kHz heartbeat layer) \cite{iec-62304}.
The current state of the art is short of the 1-minute requirement on
every key parameter by 5 to 200 times.

\subsection{On-Premises LLM Thesis}
\label{sec:intro-thesis}

The thesis of this paper is stated verbatim: on-premises repository
based LLMs provide commands to standard oncology surgical robots
based on real-time sensor data and controlled via x, y, z coordinates
to administer patient treatment, and this workflow minimizes single
robot error potential. The 1 kHz heartbeat broadcast bus and the
cumulative 4-arm tip force cap are the operational mechanisms that
turn the thesis into a measurable safety property. The 1M token
context window of Claude Code Opus 4.7 \cite{claude-opus-47,
claude-code} lets a single inference pass hold all 12 instruction
files, every kinematics and schema YAML / JSON specification, and
every per-iteration aggregate in context; on-premises inference
backends Ollama \cite{ollama} and vLLM \cite{vllm} make the same
context tractable for a deployment that keeps patient health
information on site. The 4-arm coordination heartbeat snapshot reads
as follows.

\begin{verbatim}
+-----------------------------------------------------------------+
|  1 kHz heartbeat broadcast bus  (32 byte frame, 1 ms deadline)  |
+-----------------------------------------------------------------+
|  Arm 1 hybrid u-w-p  |  Arm 2 bipolar coag  |  Arm 3 suction   |
|  cuts at 800 mm3/s   |  + irrigation        |  + collection    |
|  peak; tip <= 5.0 N  |  tip <= 5.0 N        |  tip <= 5.0 N    |
+-----------------------------------------------------------------+
|  Arm 4  0.5 T iMRI + 5-ALA fluorescence + ultrasound imaging    |
|  imaging only (no cutting force); tip <= 5.0 N (contact only)   |
+-----------------------------------------------------------------+
|  Cumulative 4-arm tip force on patient frame <= 12.0 N          |
|  E-stop 5 ms, heartbeat watchdog 3 ms, 100 us emergency park    |
+-----------------------------------------------------------------+
\end{verbatim}

An on-premises LLM that emits per-arm xyz commands at the 1 kHz tick
falls inside the FDA Software as a Medical Device (SaMD) framework
\cite{fda-samd}. The safety advantage is single-robot error
minimization through 4-arm cross-checking on the deterministic
heartbeat bus: an xyz command that would push the cumulative tip
force above the 12 N envelope is rejected at the heartbeat boundary
before it reaches the actuator, and the rejection triggers the 5 ms
E-stop with the 100 microsecond emergency arm-park. Four arms acting
as cross-checks on one another are strictly safer than a single arm
acting alone.

\subsection{Transition to the 60-Second Resection}
\label{sec:intro-transition}

The 60-second resection scenario is exercised through a 16-iteration
deterministic sweep over noise (0.01 to 0.05 mm sensor noise sigma),
force feedback gain (0.8 to 1.2), inverse-kinematics tolerance
(1e-6 to 1e-3), and heartbeat jitter (0 to 50 microseconds), all
under the fixed seed 20260510. Each iteration emits an L1 50 ms
aggregate Parquet, an L2 1 s aggregate Parquet, an L3 per-phase
aggregate Parquet, an events Parquet, and an L0 Zenodo pointer JSON
to kevinkawchak/robotic-surgeries/2030-gbm-1min/outputs/iterations/.
Two on-premises LLM tournaments run against the iteration output:
the default 4-entity robot-vs-robot tournament at
2030-gbm-1min/outputs/comparison/ (6 rounds) and the mixed 2 robot
plus 2 human tournament at
2030-gbm-1min/outputs/comparison_robot_vs_human/ (6 rounds). The
headline result that the paper builds toward is that the robot mean
composite score is 88.53 versus human mean composite 70.35, with the
robot winning all 4 robot-vs-human pairings at confidence 0.955 to
1.000; the structural-time-dimension caveat (1-minute robot vs
1-hour human baseline) is preserved in every LLM judge rationale
\cite{kawchak_2026_20113157, repo-robotic-surgeries}. The on-screen
Table of Contents that opens page 2 lists the remaining sections.
